\section{Despliegue y observación}

\subsection{Monitoreo de métricas inference-time}
Chen ve el inference-time pipeline como una «caja negra»: tras recibir el prompt/block de entrada, pasa sucesivamente por sampling, search y trace, y cada etapa debe generar metrics. Las metrics típicas incluyen sampling depth, branching ratio, verifier score distribution, trace iteration count y latency. Se recomienda graficar estas curvas en el dashboard agrupadas cada 5 minutos; en un sistema real, en cuanto el sampling depth caiga de manera repentina o el trace iteration aumente bruscamente, debe activarse una investigación.

\begin{knowledgebox}{Cadena de herramientas de observabilidad}
1. Sampling log: registra prompt + generated candidates + token count; 2. Search log: registra beam/tree status (depth, score, branch id); 3. Trace log: indica en cada ronda de self-correction si hay un error pattern nombrado por humanos; 4. Alert rules: configura un collector que emita una alerta cuando la diversidad de sample sea inferior al baseline o el trace drift supere el threshold.
\end{knowledgebox}

\subsection{Budget SLO y alertas}
Al desplegar el inference-time compute pipeline, es común usar una tabla de SLO que enumere los límites superiores de token budget, search depth e iteration depth, junto con alertas del tipo «cuando average token per query > budget × 1.1» o «search branching > 30%». Chen recomienda agregar una etapa de «double-check» en el SLO, y usar los dos score del stepwise verifier para asegurar que una high-stakes query no falle debido a un sampling drift.

\begin{table}[H]
\centering
\caption{Ejemplo de SLO de inference-time}
\begin{tabular}{p{0.32\textwidth}p{0.32\textwidth}p{0.32\textwidth}}
\toprule
Metric & Target & Alert condition \\
\midrule
Token budget per query & $\leq 400$ tokens & 5-minute moving average > 440 tokens \\
Sampling diversity & 80\% unique candidates & 3 consecutive minutes diversity < 60\% \\
Iteration depth & $\leq 2$ feedback rounds & iteration count > 3 for top 1\% queries \\
Verifier score margin & $\geq 0.15$ & margin drop > 0.05 for 100 queries \\
\bottomrule
\end{tabular}
\end{table}

\begin{warningbox}{Recomendaciones para el ajuste de alertas}
No permitas que cada trace drift dispare una alerta, pues eso provocaría alarm fatigue. La experiencia de Chen es fijar el umbral de falsos positivos un poco más alto, y programar una investigación manual solo cuando exista un problema real, apoyándose en la función de «replay» para que el equipo revise el sampling log.
\end{warningbox}

\subsection{Resumen de la sección}
Al desplegar el inference-time pipeline es necesario que observabilidad, SLO, alertas y log replay trabajen de manera coordinada, para poder mantener accuracy y reliability en la capa de production, y al mismo tiempo mantener el budget dentro de un rango controlable.

